% ----------------------------------------------------------
% Considerações Finais
% ----------------------------------------------------------
\chapter{Considerações Finais}

O desenvolvimento do \textit{SyncCarreira} representou não apenas um desafio técnico de engenharia de \textit{software}, mas, sobretudo, uma oportunidade de aplicar a tecnologia como ferramenta de transformação social. Ao longo deste projeto, foi possível estruturar uma solução capaz de democratizar o acesso à orientação profissional, utilizando uma arquitetura robusta e escalável, pensada para atender às necessidades específicas de jovens em contextos de vulnerabilidade.

\section{Dificuldades, Escolhas e Descartes}

A construção do \textit{SyncCarreira} foi um processo de aprendizado contínuo, onde a tradução do negócio para a tecnologia foi o desafio central.

\begin{itemize}
    \item \textbf{Dificuldades:} A principal barreira foi o alinhamento de expectativas e a tradução das necessidades da instituição parceira para requisitos técnicos. O processo de orientação vocacional, por ser inicialmente abstrato, exigiu um esforço extra de comunicação e modelagem para que as regras de negócio fossem compreendidas pela equipe de desenvolvimento e convertidas em funcionalidades concretas e tangíveis no sistema.

    \item \textbf{Escolhas:} Optou-se pela utilização do MySQL como banco de dados relacional pela sua robustez e confiabilidade na persistência de dados. A escolha do React no \textit{front-end} foi estratégica para garantir que a interface fosse tecnicamente consistente e fácil de manter a longo prazo, permitindo uma estrutura de componentes clara e modular.

    \item \textbf{Descartes:} Durante o refinamento do projeto, descartou-se a implementação de funcionalidades que prometiam automatizações complexas de Inteligência Artificial (IA), optando-se por priorizar a clareza do fluxo de navegação e a curadoria de conteúdo pedagógico. Esse descarte foi fundamental para evitar a criação de um sistema excessivamente complexo e garantir que a ferramenta entregasse valor imediato aos jovens atendidos pela ONG, sem dispersar o foco da equipe de desenvolvimento em tecnologias que não eram vitais para o MVP.
\end{itemize}
